% Front matter: title verso (edition notice, Sec 10) and About This Book.
% All boilerplate newly written for Initech (DESIGN-SYSTEM.md Sec 10).
\begin{TitleVerso}
{\imLegalFont
\textbf{INITECH SYSTEMS CORPORATION PROVIDES THIS PUBLICATION ``AS IS''
WITHOUT WARRANTY OF ANY KIND, EITHER EXPRESS OR IMPLIED, INCLUDING, BUT NOT
LIMITED TO, THE IMPLIED WARRANTIES OF MERCHANTABILITY AND FITNESS FOR A
PARTICULAR PURPOSE.}  Some states do not allow the exclusion of implied
warranties, so the above exclusion may not apply to you.\par
This publication may contain technical inaccuracies or typographical errors.
The information in it is revised from time to time; revisions are issued as
new editions or technical newsletters.  Initech Systems Corporation may change
the products and programs described in this publication at any time without
notice.  A reference to an Initech product or program does not imply that only
that product or program may be used.\par
Requests for copies of Initech publications and for technical information
about Initech products should be made to your authorized Initech dealer or
your Initech representative.\par}

{\footnotesize
A Reader's Comment Form is provided at the back of this publication.  If the
form has been removed, address comments to: \InitechAddressLine.
\draft{D-12: confirm the address belongs to no real business before print}
Initech may use or distribute the information you supply in any way it
believes appropriate without incurring any obligation to you.\par}

{\imLegalFont
InitechDOS, FLAIR, InitechBase and the Initech logo are trademarks of
Initech Systems Corporation.  Turbo Initech is a registered trademark of
Initech Systems Corporation.  Intel, 80386 and 80486 are trademarks of Intel
Corporation.  All other product names are trademarks of their respective
holders.\par}

{\footnotesize
This book was typeset by Initech Publications with LuaTeX (TeX Live) in TeX
Gyre Schola, TeX Gyre Heros and TeX Gyre Cursor, metric-compatible
replacements for New Century Schoolbook, Helvetica and Courier.  The screen
reproductions are the text of the InitechDOS display, read back from
screendumps of the shipped system running under the QEMU emulator.\par}
\end{TitleVerso}

\FrontChapter{About This Book}
% Gates: none (orientation text); the conventions list is checked against
% pubs/class/initech-syntax.sty.
This book shows you how to start InitechDOS, how to give it commands, and how
to look after the files you keep on your disks.  It is written for people who
have not used a personal computer operating system before, and for
experienced users who are new to InitechDOS.  You do not need to know
anything about programming.

\section{Read This First}
Before you begin, make sure your computer has been set up as described in
the instructions that came with it, and that you have the InitechDOS
diskettes from the package.  Keep this book beside your computer while you
work: Chapter~2 is written as a series of short exercises that you carry
out as you read.

\section{How This Book Is Organized}
\SideHead{Chapter 1}Introduction to InitechDOS explains what an operating
system does, names the parts of InitechDOS, and introduces drives, files,
directories and commands.

\SideHead{Chapter 2}Getting Started takes you from turning on the computer
to listing, displaying, copying and deleting files, and to working with
directories.

\SideHead{Appendix A}Messages lists the messages InitechDOS displays.

An index and a Reader's Comment Form are at the back of the book.

\section{Related Publications}
\SideHead{Reference}\doctitle{InitechDOS Version 3.30 Reference}
(1024-3300-11) describes every command and configuration option in detail.

\SideHead{Quick Reference}\doctitle{InitechDOS Quick Reference}
(1024-3300-20) summarizes the commands on a single card.
\draft{the Reference and Quick Reference are planned, not yet written}

\section{Conventions Used in This Book}
\SideHead{What you type}Text that you type is shown in this typeface:
\type{DIR}.  You may type commands in capital or small letters; InitechDOS
treats them the same.  Unless the text tells you otherwise, finish each
command by pressing \key{Enter}.

\SideHead{Keys}The names of keys are shown in small capitals, for example
\key{Enter}, \key{Backspace} and \key{Ctrl}.  \key{Ctrl-C} means: hold down
\key{Ctrl}, press \key{C}, and release both keys.

\SideHead{Command syntax}When this book shows the form of a command, words in
capital letters, such as \kw{COPY}, are typed exactly as shown.  Words in
\ph{italics} stand for information you supply, such as a file name.  Items
in square brackets \opt{\ph{like this}} are optional; do not type the
brackets.  Three periods \rep{} mean that the item before them may be
repeated.

\SideHead{Screens}Figures that show the display are reproductions of the
InitechDOS screen exactly as it appeared when the example was carried out.
